% Generated draft: requires mathematical and source review.
\documentclass[11pt]{article}
\usepackage[margin=25mm]{geometry}
\usepackage{fontspec}
\setmainfont{Latin Modern Roman}
\usepackage{amsmath,amssymb,mathtools,unicode-math}
\setmathfont{Latin Modern Math}
\usepackage{enumitem,needspace,xurl,hyperref}
\setlength{\parindent}{0pt}
\setlength{\parskip}{5pt}
\setlength{\emergencystretch}{3em}
\newcommand{\N}{\mathbb{N}}
\newcommand{\Z}{\mathbb{Z}}
\newcommand{\Q}{\mathbb{Q}}
\newcommand{\R}{\mathbb{R}}
\newcommand{\C}{\mathbb{C}}
\newcommand{\F}{\mathbb{F}}
\newcommand{\A}{\mathbb{A}}
\newcommand{\PP}{\mathbb{P}}
\newcommand{\E}{\mathbb{E}}
\newcommand{\eps}{\varepsilon}
\newcommand{\dd}{\,\mathrm{d}}
\begin{document}
\section*{UnsolvedMath ID: 30000786}
\textbf{Source problem code:} catalogue entry
\subsection{Definitions and mathematical statement}
For a finite simple graph \(G\) on \(N\) vertices, define
\[
 p(G)=\frac{|E(G)|}{\binom N2},\qquad
 t_r(G)=\frac{\#\{S\subseteq V(G):|S|=r,\ G[S]\cong K_r\}}{\binom Nr}.
\]
Thus each unordered vertex set is counted once. With the usual asymptotic interpretation, the minimum triangle-density problem asks to determine
\[
 f_3(p)=\inf\left\{\liminf_{j\to\infty}t_3(G_j):
 |V(G_j)|\to\infty,\ p(G_j)\to p\right\}
 \qquad\left(\frac12<p\leq1\right).
\]
Here \(1/2\) is the ordinary graph triangle-free critical edge density. The record does not explicitly choose an asymptotic versus finite-size formulation or a density normalization; the displayed convention makes the standard asymptotic request precise.

For a finite simple three-uniform hypergraph \(H\), edges are three-element subsets of its vertex set. Let \(K_4^{(3)}\) have all four triples on four vertices and put
\[
 \operatorname{ex}(N,K_4^{(3)})
 =\max\{|E(H)|:|V(H)|=N,\ H\text{ contains no }K_4^{(3)}\},
\]
\[
 \pi(K_4^{(3)})
 =\lim_{N\to\infty}\frac{\operatorname{ex}(N,K_4^{(3)})}{\binom N3}.
\]
Containment here is as a subhypergraph, not an additional induced-forbidden condition. The second explicit request is to determine this Tur\'an density.

The source also requests analogous \(K_r\) questions for \(r>3\). In the ordinary graph setting the natural analogue is the corresponding function \(f_r(p)\), obtained by replacing \(t_3\) with \(t_r\), above the graph \(K_r\)-free threshold \(1-1/(r-1)\). In a three-uniform setting, analogous forbidden complete hypergraphs would be denoted \(K_r^{(3)}\). The supplied shorthand \(K_r\) does not identify which of these scopes, or which varying-uniformity variant, is intended, so these interpretations are recorded rather than conflated.

The supplied assessment dated 2026-08-21 calls the ordinary triangle question classical, reports important partial flag-algebra results for the three-uniform question, and says the general \(K_4^{(3)}\) Tur\'an-density problem remains open. It does not supply an exact value or a complete theorem for the higher-clique request.
\subsection{Short English statement}
The record combines a graph clique-density minimization problem at a prescribed edge density with a hypergraph problem maximizing the density while avoiding a complete four-vertex three-uniform hypergraph. It also asks for higher-clique analogues, whose graph or hypergraph scope needs clarification.
\subsection{Sources}
\begin{itemize}
\item Local UnsolvedMath registry, permanent ID \(30000786\), exact source code \texttt{catalogue entry}, proposed year 2007; no individual proposer is supplied.
\item \emph{Complexity Theory (2008)}, Oberwolfach Reports, DOI 10.4171/owr/2007/31, pp. 1793--1864. \url{https://doi.org/10.4171/owr/2007/31}.
\item Supplied reference: Razborov, ``forbidden 4-vertex configurations,'' \url{https://doi.org/10.1137/090747476}.
\item The \texttt{corrected\_verified} note says the full density request was retained while uniformity and parameter notation were restored. The partially solved label and assessment checked 2026-08-21 are dated supplied literature triage requiring independent verification of status and theorem scope.
\end{itemize}
\end{document}
